\documentclass{article}
\usepackage[T1]{fontenc}
\usepackage{lmodern}
\usepackage{graphicx}
\usepackage{amsmath,amssymb}
\usepackage[a4paper,margin=2cm]{geometry}
\usepackage[absolute,overlay]{textpos}
\setlength{\TPHorizModule}{1mm}
\setlength{\TPVertModule}{1mm}
\usepackage[indonesian]{babel}
\usepackage{hyperref}
\setlength{\emergencystretch}{2em}
% unit: Halaman 17

\begin{document}

% === Halaman 17 (python/native) ===

(b) Enkripsi  (oleh Bob)

Misalkan Bob ingin mengirim plainteks ‘HALO’ kepada Alice. 

Misalkan A = 00, B = 01, …, Z = 25, maka pesan m dikodekan ke dalam integer  

adalah 

m = 07001114

Bob memecah m menjadi blok yang lebih kecil, misalkan m dipecah menjadi blok-

blok sepanjang 4 angka:

m1 = 0700

dan m2 = 1114

Nilai-nilai mi ini masih terletak di dalam selang [0, 2273 – 1] agar transformasi 

menjadi satu-ke-satu. 

17

\end{document}
